\documentclass[10pt,a4paper]{amsart}
\usepackage{amsmath,amssymb,amsthm,graphicx}
% Retained preamble for the Stacks Project chapter snapshot in this directory.
% The chapter's own preamble.tex is a copy of the shared Stacks Project preamble; the block
% below is sliced MECHANICALLY, line for line, from it. Nothing was added, removed or
% redefined; the wrapper class, the xr-hyper/hyperref wiring and the externaldocument list
% are omitted because the assembler supplies the document class.
%
% lines taken from preamble.tex: 3, 5, 9, 11, 12, 156, 157, 158, ... (51 lines total)
%
% package coverage restored from the chapter's own preamble (2026-10-13): xy, xy-2cell

\usepackage[all]{xy}
\xyoption{2cell}
\UseAllTwocells

\newenvironment{reference}{\comment}{\endcomment}
\newenvironment{slogan}{\comment}{\endcomment}
\newenvironment{history}{\comment}{\endcomment}
\theoremstyle{plain}
\newtheorem{theorem}[subsection]{Theorem}
\newtheorem{proposition}[subsection]{Proposition}
\newtheorem{lemma}[subsection]{Lemma}
\theoremstyle{definition}
\newtheorem{definition}[subsection]{Definition}
\newtheorem{example}[subsection]{Example}
\newtheorem{exercise}[subsection]{Exercise}
\newtheorem{situation}[subsection]{Situation}
\theoremstyle{remark}
\newtheorem{remark}[subsection]{Remark}
\newtheorem{remarks}[subsection]{Remarks}
\numberwithin{equation}{subsection}
\def\lim{\mathop{\mathrm{lim}}\nolimits}
\def\colim{\mathop{\mathrm{colim}}\nolimits}
\def\Spec{\mathop{\mathrm{Spec}}}
\def\Hom{\mathop{\mathrm{Hom}}\nolimits}
\def\Ext{\mathop{\mathrm{Ext}}\nolimits}
\def\SheafHom{\mathop{\mathcal{H}\!\mathit{om}}\nolimits}
\def\SheafExt{\mathop{\mathcal{E}\!\mathit{xt}}\nolimits}
\def\Sch{\mathit{Sch}}
\def\Mor{\mathop{\mathrm{Mor}}\nolimits}
\def\Ob{\mathop{\mathrm{Ob}}\nolimits}
\def\Sh{\mathop{\mathit{Sh}}\nolimits}
\def\NL{\mathop{N\!L}\nolimits}
\def\CH{\mathop{\mathrm{CH}}\nolimits}
\def\proetale{{pro\text{-}\acute{e}tale}}
\def\etale{{\acute{e}tale}}
\def\QCoh{\mathit{QCoh}}
\def\Ker{\mathop{\mathrm{Ker}}}
\def\Im{\mathop{\mathrm{Im}}}
\def\Coker{\mathop{\mathrm{Coker}}}
\def\Coim{\mathop{\mathrm{Coim}}}
\DeclareMathSymbol{\boxtimes}{\mathbin}{AMSa}{"02}
\def\QCohstack{\mathcal{QC}\!\mathit{oh}}
\def\Cohstack{\mathcal{C}\!\mathit{oh}}
\def\Spacesstack{\mathcal{S}\!\mathit{paces}}
\def\Quotfunctor{\mathrm{Quot}}
\def\Hilbfunctor{\mathrm{Hilb}}
\def\Curvesstack{\mathcal{C}\!\mathit{urves}}
\def\Polarizedstack{\mathcal{P}\!\mathit{olarized}}
\def\Complexesstack{\mathcal{C}\!\mathit{omplexes}}
\def\Pic{\mathop{\mathrm{Pic}}\nolimits}
\def\Picardstack{\mathcal{P}\!\mathit{ic}}
\def\Picardfunctor{\mathrm{Pic}}
\def\Deformationcategory{\mathcal{D}\!\mathit{ef}}
\usepackage{cleveref}
\crefname{originalproposition}{original proposition}{original propositions}
\Crefname{originalproposition}{Original proposition}{Original propositions}
\makeatletter
\newlabel{example-universal-quasi-finite-syntomic}{{0FK8}{source}{}{sourcepointer-1}{}}
\newlabel{example-universal-quasi-finite-syntomic@cref}{{[originalproposition][0][]0FK8}{[1][0][]source}{}{}{}}
\makeatother
\crefname{originalproposition}{original proposition}{original propositions}
\Crefname{originalproposition}{Original proposition}{Original propositions}
\makeatletter
\newlabel{lemma-different-quasi-finite-complete-intersection}{{0BWH}{source}{}{sourcepointer-2}{}}
\newlabel{lemma-different-quasi-finite-complete-intersection@cref}{{[originalproposition][0][]0BWH}{[1][0][]source}{}{}{}}
\makeatother
\crefname{originalproposition}{original proposition}{original propositions}
\Crefname{originalproposition}{Original proposition}{Original propositions}
\makeatletter
\newlabel{lemma-different-ramification}{{0BW9}{source}{}{sourcepointer-3}{}}
\newlabel{lemma-different-ramification@cref}{{[originalproposition][0][]0BW9}{[1][0][]source}{}{}{}}
\makeatother
\crefname{originalproposition}{original proposition}{original propositions}
\Crefname{originalproposition}{Original proposition}{Original propositions}
\makeatletter
\newlabel{lemma-locally-comes-from-universal}{{0FKA}{source}{}{sourcepointer-4}{}}
\newlabel{lemma-locally-comes-from-universal@cref}{{[originalproposition][0][]0FKA}{[1][0][]source}{}{}{}}
\makeatother
\crefname{originalproposition}{original proposition}{original propositions}
\Crefname{originalproposition}{Original proposition}{Original propositions}
\makeatletter
\newlabel{lemma-syntomic-quasi-finite}{{0BWE}{source}{}{sourcepointer-5}{}}
\newlabel{lemma-syntomic-quasi-finite@cref}{{[originalproposition][0][]0BWE}{[1][0][]source}{}{}{}}
\makeatother
\crefname{originalproposition}{original proposition}{original propositions}
\Crefname{originalproposition}{Original proposition}{Original propositions}
\makeatletter
\newlabel{lemma-universal-quasi-finite-syntomic-etale}{{0FK9}{source}{}{sourcepointer-6}{}}
\newlabel{lemma-universal-quasi-finite-syntomic-etale@cref}{{[originalproposition][0][]0FK9}{[1][0][]source}{}{}{}}
\makeatother
\crefname{originalproposition}{original proposition}{original propositions}
\Crefname{originalproposition}{Original proposition}{Original propositions}
\makeatletter
\newlabel{limits-lemma-descend-finite-presentation}{{01ZM}{source}{}{sourcepointer-7}{}}
\newlabel{limits-lemma-descend-finite-presentation@cref}{{[originalproposition][0][]01ZM}{[1][0][]source}{}{}{}}
\makeatother
\crefname{originalproposition}{original proposition}{original propositions}
\Crefname{originalproposition}{Original proposition}{Original propositions}
\makeatletter
\newlabel{limits-lemma-descend-quasi-finite}{{094M}{source}{}{sourcepointer-8}{}}
\newlabel{limits-lemma-descend-quasi-finite@cref}{{[originalproposition][0][]094M}{[1][0][]source}{}{}{}}
\makeatother
\crefname{originalproposition}{original proposition}{original propositions}
\Crefname{originalproposition}{Original proposition}{Original propositions}
\makeatletter
\newlabel{limits-lemma-descend-syntomic}{{0C3L}{source}{}{sourcepointer-9}{}}
\newlabel{limits-lemma-descend-syntomic@cref}{{[originalproposition][0][]0C3L}{[1][0][]source}{}{}{}}
\makeatother
\crefname{originalproposition}{original proposition}{original propositions}
\Crefname{originalproposition}{Original proposition}{Original propositions}
\makeatletter
\newlabel{more-algebra-lemma-module-over-fibre-product}{{07RU}{source}{}{sourcepointer-10}{}}
\newlabel{more-algebra-lemma-module-over-fibre-product@cref}{{[originalproposition][0][]07RU}{[1][0][]source}{}{}{}}
\makeatother
\crefname{originalproposition}{original proposition}{original propositions}
\Crefname{originalproposition}{Original proposition}{Original propositions}
\makeatletter
\newlabel{more-algebra-lemma-properties-algebras-over-fibre-product}{{08KQ}{source}{}{sourcepointer-11}{}}
\newlabel{more-algebra-lemma-properties-algebras-over-fibre-product@cref}{{[originalproposition][0][]08KQ}{[1][0][]source}{}{}{}}
\makeatother
\crefname{originalproposition}{original proposition}{original propositions}
\Crefname{originalproposition}{Original proposition}{Original propositions}
\makeatletter
\newlabel{properties-lemma-standard-open-two-affines}{{0H9B}{source}{}{sourcepointer-12}{}}
\newlabel{properties-lemma-standard-open-two-affines@cref}{{[originalproposition][0][]0H9B}{[1][0][]source}{}{}{}}
\makeatother
\crefname{originalproposition}{original proposition}{original propositions}
\Crefname{originalproposition}{Original proposition}{Original propositions}
\makeatletter
\newlabel{remark-local-description-delta}{{0FKC}{source}{}{sourcepointer-13}{}}
\newlabel{remark-local-description-delta@cref}{{[originalproposition][0][]0FKC}{[1][0][]source}{}{}{}}
\makeatother
\title{Stacks, Discriminants: proposition-tate-map}
\author{The Stacks Project Authors}
\date{stacks-project:discriminant.tex:local-snapshot-20261012; GNU FDL 1.2 or later}
\begin{document}
\maketitle
\noindent\textit{Editorial note.} Statement span is lines 2933-2946 of sources/stacks-discriminant/pinned-original-source.tex; proof span is lines 2948-3163, the complete author proof (214 body lines).\par\medskip
\noindent\textit{Editorial note.} Navigation stubs: 13 cited-but-not-included labels are supplied as non-load-bearing pointers with their pinned tags where the snapshot tag index has them.\par\medskip
\noindent\textit{Editorial note.} extra\_packages is empty on purpose: the retained preamble already loads xy with options and a second declaration clashes in the assembler wrapper.\par\medskip
\noindent\textit{Editorial note.} Text is retained verbatim, including the author's own typos and placeholders.\par\medskip
\noindent\textit{Editorial note.} Queue index 0 of queue\_discriminant.json (pb=9070, stmt 2933-2946, proof 2948-3163).\par\medskip
\noindent\textit{Editorial note.} Cross-chapter same-label warning (NOT an occupancy conflict, Chief ruling 2026-10-13 evening): label 'proposition-tate-map' also exists in stacks-snapshot-20261011 — a different proposition in a different chapter file.\par\medskip
\noindent\textit{Allowed standard prerequisites (editorial).}\begin{itemize}
\item Explicit prerequisite boundary. The cited inputs are prerequisites; the argument in the retained span is not replaced by any cited result.
\end{itemize}\medskip
\ifdefined\hypertarget\hypertarget{sourcepointer-1}{}\fi
\noindent\textit{Source reference.} Original proposition 0FK8: sources/stacks-discriminant/pinned-original-source.tex, source lines 2168--2185. Cited by the retained statement or proof but not part of the collected spans; supplied only as a navigation pointer. The argument does not depend on its content.\par\medskip
\ifdefined\hypertarget\hypertarget{sourcepointer-2}{}\fi
\noindent\textit{Source reference.} Original proposition 0BWH: sources/stacks-discriminant/pinned-original-source.tex, source lines 2727--2734. Cited by the retained statement or proof but not part of the collected spans; supplied only as a navigation pointer. The argument does not depend on its content.\par\medskip
\ifdefined\hypertarget\hypertarget{sourcepointer-3}{}\fi
\noindent\textit{Source reference.} Original proposition 0BW9: sources/stacks-discriminant/pinned-original-source.tex, source lines 1907--1912. Cited by the retained statement or proof but not part of the collected spans; supplied only as a navigation pointer. The argument does not depend on its content.\par\medskip
\ifdefined\hypertarget\hypertarget{sourcepointer-4}{}\fi
\noindent\textit{Source reference.} Original proposition 0FKA: sources/stacks-discriminant/pinned-original-source.tex, source lines 2225--2242. Cited by the retained statement or proof but not part of the collected spans; supplied only as a navigation pointer. The argument does not depend on its content.\par\medskip
\ifdefined\hypertarget\hypertarget{sourcepointer-5}{}\fi
\noindent\textit{Source reference.} Original proposition 0BWE: sources/stacks-discriminant/pinned-original-source.tex, source lines 2005--2026. Cited by the retained statement or proof but not part of the collected spans; supplied only as a navigation pointer. The argument does not depend on its content.\par\medskip
\ifdefined\hypertarget\hypertarget{sourcepointer-6}{}\fi
\noindent\textit{Source reference.} Original proposition 0FK9: sources/stacks-discriminant/pinned-original-source.tex, source lines 2188--2195. Cited by the retained statement or proof but not part of the collected spans; supplied only as a navigation pointer. The argument does not depend on its content.\par\medskip
\ifdefined\hypertarget\hypertarget{sourcepointer-7}{}\fi
\noindent\textit{Source reference.} Original proposition 01ZM: sources/stacks-discriminant/pinned-original-source.tex, source lines 3163--3163. Cited by the retained statement or proof but not part of the collected spans; supplied only as a navigation pointer. The argument does not depend on its content.\par\medskip
\ifdefined\hypertarget\hypertarget{sourcepointer-8}{}\fi
\noindent\textit{Source reference.} Original proposition 094M: sources/stacks-discriminant/pinned-original-source.tex, source lines 3163--3163. Cited by the retained statement or proof but not part of the collected spans; supplied only as a navigation pointer. The argument does not depend on its content.\par\medskip
\ifdefined\hypertarget\hypertarget{sourcepointer-9}{}\fi
\noindent\textit{Source reference.} Original proposition 0C3L: sources/stacks-discriminant/pinned-original-source.tex, source lines 3163--3163. Cited by the retained statement or proof but not part of the collected spans; supplied only as a navigation pointer. The argument does not depend on its content.\par\medskip
\ifdefined\hypertarget\hypertarget{sourcepointer-10}{}\fi
\noindent\textit{Source reference.} Original proposition 07RU: sources/stacks-discriminant/pinned-original-source.tex, source lines 3163--3163. Cited by the retained statement or proof but not part of the collected spans; supplied only as a navigation pointer. The argument does not depend on its content.\par\medskip
\ifdefined\hypertarget\hypertarget{sourcepointer-11}{}\fi
\noindent\textit{Source reference.} Original proposition 08KQ: sources/stacks-discriminant/pinned-original-source.tex, source lines 3163--3163. Cited by the retained statement or proof but not part of the collected spans; supplied only as a navigation pointer. The argument does not depend on its content.\par\medskip
\ifdefined\hypertarget\hypertarget{sourcepointer-12}{}\fi
\noindent\textit{Source reference.} Original proposition 0H9B: sources/stacks-discriminant/pinned-original-source.tex, source lines 3163--3163. Cited by the retained statement or proof but not part of the collected spans; supplied only as a navigation pointer. The argument does not depend on its content.\par\medskip
\ifdefined\hypertarget\hypertarget{sourcepointer-13}{}\fi
\noindent\textit{Source reference.} Original proposition 0FKC: sources/stacks-discriminant/pinned-original-source.tex, source lines 2867--2898. Cited by the retained statement or proof but not part of the collected spans; supplied only as a navigation pointer. The argument does not depend on its content.\par\medskip
\begin{proposition}
\label{proposition-tate-map}
There exists a unique rule that to every locally quasi-finite syntomic
morphism of locally Noetherian schemes $Y \to X$ assigns an isomorphism
$$
c_{Y/X} : \det(\NL_{Y/X}) \longrightarrow \omega_{Y/X}
$$
satisfying the following two properties
\begin{enumerate}
\item the section $\delta(\NL_{Y/X})$ is mapped to $\tau_{Y/X}$, and
\item the rule is compatible with restriction to opens and with
base change.
\end{enumerate}
\end{proposition}

\begin{proof}
Let us reformulate the statement of the proposition. Consider the category
$\mathcal{C}$ whose objects, denoted $Y/X$, are locally quasi-finite syntomic
morphism $Y \to X$ of locally Noetherian schemes and whose morphisms
$b/a : Y'/X' \to Y/X$ are commutative diagrams
$$
\xymatrix{
Y' \ar[d] \ar[r]_b & Y \ar[d] \\
X' \ar[r]^a & X
}
$$
which induce an isomorphism of $Y'$ with an open subscheme of
$X' \times_X Y$. The proposition means that for every object
$Y/X$ of $\mathcal{C}$ we have an isomorphism
$c_{Y/X} : \det(\NL_{Y/X}) \to \omega_{Y/X}$
with $c_{Y/X}(\delta(\NL_{Y/X})) = \tau_{Y/X}$
and for every morphism $b/a : Y'/X' \to Y/X$ of $\mathcal{C}$ we have
$b^*c_{Y/X} = c_{Y'/X'}$ via the identifications
$b^*\det(\NL_{Y/X}) = \det(\NL_{Y'/X'})$ and
$b^*\omega_{Y/X} = \omega_{Y'/X'}$ described above.

\medskip\noindent
Given $Y/X$ in $\mathcal{C}$ and $y \in Y$ we can find
an affine opens $V \subset Y$ and $U \subset X$ with
$y \in V$ and $f(V) \subset U$
such that there exists some isomorphism
$$
\det(\NL_{Y/X})|_V \longrightarrow \omega_{Y/X}|_V
$$
mapping $\delta(\NL_{Y/X})|_V$ to $\tau_{Y/X}|_V$. This follows
from picking affine opens as in
Lemma \ref{lemma-syntomic-quasi-finite} part (5), the affine
local description of $\delta(\NL_{Y/X})$ in
Remark \ref{remark-local-description-delta}, and
Lemma \ref{lemma-different-quasi-finite-complete-intersection}.
If the annihilator of the section $\tau_{Y/X}$ is zero, then
these local maps are unique and automatically glue. Hence if the annihilator
of $\tau_{Y/X}$ is zero, then there is a unique isomorphism
$c_{Y/X} : \det(\NL_{Y/X}) \to \omega_{Y/X}$ with
$c_{Y/X}(\delta(\NL_{Y/X})) = \tau_{Y/X}$.
If $b/a : Y'/X' \to Y/X$ is a morphism of $\mathcal{C}$
and the annihilator of $\tau_{Y'/X'}$ is zero as well,
then $b^*c_{Y/X}$ is the unique isomorphism
$c_{Y'/X'} : \det(\NL_{Y'/X'}) \to \omega_{Y'/X'}$ with
$c_{Y'/X'}(\delta(\NL_{Y'/X'})) = \tau_{Y'/X'}$.
This follows formally from the fact that
$b^*\delta(\NL_{Y/X}) = \delta(\NL_{Y'/X'})$ and
$b^*\tau_{Y/X} = \tau_{Y'/X'}$.

\medskip\noindent
We can summarize the results of the previous paragraph as follows.
Let $\mathcal{C}_{nice} \subset \mathcal{C}$ denote the
full subcategory of $Y/X$ such that the annihilator of
$\tau_{Y/X}$ is zero. Then we have solved the problem
on $\mathcal{C}_{nice}$. For $Y/X$ in $\mathcal{C}_{nice}$
we continue to denote $c_{Y/X}$ the solution we've just found.

\medskip\noindent
Consider morphisms
$$
Y_1/X_1 \xleftarrow{b_1/a_1} Y/X \xrightarrow{b_2/a_2} Y_2/X_2
$$
in $\mathcal{C}$ such that $Y_1/X_1$ and $Y_2/X_2$ are objects
of $\mathcal{C}_{nice}$. {\bf Claim.} $b_1^*c_{Y_1/X_1} = b_2^*c_{Y_2/X_2}$.
We will first show that the claim implies the proposition
and then we will prove the claim.

\medskip\noindent
Let $d, n \geq 1$ and consider the locally
quasi-finite syntomic morphism $Y_{n, d} \to X_{n, d}$
constructed in Example \ref{example-universal-quasi-finite-syntomic}.
Then $Y_{n, d}$ is an irreducible regular scheme and the
morphism $Y_{n, d} \to X_{n, d}$ is locally quasi-finite syntomic
and \'etale over a dense open, see
Lemma \ref{lemma-universal-quasi-finite-syntomic-etale}.
Thus $\tau_{Y_{n, d}/X_{n, d}}$ is nonzero for example by
Lemma \ref{lemma-different-ramification}. Now a nonzero section
of an invertible module over an irreducible regular scheme
has vanishing annihilator. Thus
$Y_{n, d}/X_{n, d}$ is an object of $\mathcal{C}_{nice}$.

\medskip\noindent
Let $Y/X$ be an arbitrary object of $\mathcal{C}$. Let $y \in Y$.
By Lemma \ref{lemma-locally-comes-from-universal} we can find
$n, d \geq 1$ and morphisms
$$
Y/X \leftarrow V/U \xrightarrow{b/a} Y_{n, d}/X_{n, d}
$$
of $\mathcal{C}$ such that $V \subset Y$ and $U \subset X$ are open.
Thus we can pullback the canonical morphism $c_{Y_{n, d}/X_{n, d}}$
constructed above by $b$ to $V$. The claim guarantees these local
isomorphisms glue! Thus we get a well defined global isomorphism
$c_{Y/X} : \det(\NL_{Y/X}) \to \omega_{Y/X}$ with
$c_{Y/X}(\delta(\NL_{Y/X})) = \tau_{Y/X}$.
If $b/a : Y'/X' \to Y/X$ is a morphism of $\mathcal{C}$, then
the claim also implies that the similarly constructed map
$c_{Y'/X'}$ is the pullback by $b$ of the locally constructed
map $c_{Y/X}$. Thus it remains to prove the claim.

\medskip\noindent
In the rest of the proof we prove the claim. We may pick a point
$y \in Y$ and prove the maps agree in an open neighbourhood of $y$.
Thus we may replace $Y_1$, $Y_2$ by open neighbourhoods of the
image of $y$ in $Y_1$ and $Y_2$. Therefore we may assume
$Y, X, Y_1, X_1, Y_2, X_2$ are affine, say they are the spectra
of rings $B, A, B_1, A_1, B_2, A_2$. Picture
$$
\xymatrix{
B_1 \ar[r] & B & B_2 \ar[l] \\
A_1 \ar[u] \ar[r] & A \ar[u] & A_2 \ar[l] \ar[u]
}
$$
By assumption the spectrum of $B$ is an affine open of
both the spectrum of $A \otimes_{A_1} B_1$ and $A \otimes_{A_2} B_2$.
Shrinking more we may assume there exist elements
$g_i \in A \otimes_{A_i} B_i$ such that our maps give isomorphisms
$(A \otimes_{A_i} B_i)_{g_i} = B$, see
Properties, Lemma \ref{properties-lemma-standard-open-two-affines}.
Let $x_\alpha$, $y_\beta$ be a sufficiently large collection of
variables such that we may choose surjections
$$
A'_1 = A_1[x_\alpha] \to A
\quad\text{and}\quad
A'_2 = A_2[y_\beta] \to A
$$
of $A_1$ and $A_2$-algebras. Then we can choose lifts
$h_i \in A'_i \otimes_{A_i} B_i$ of $g_i \in A \otimes_{A_i} B_i$
and we consider the diagram
$$
\xymatrix{
(A'_1 \otimes_{A_1} B_1)_{h_1} \ar[r] & B &
(A'_2 \otimes_{A_1} B_2)_{h_2} \ar[l] \\
A'_1 \ar[u] \ar[r] & A \ar[u] & A'_2 \ar[l] \ar[u]
}
$$
By construction the two squares are cocartesian. Next, we consider the
ring map
$$
A' = A'_1 \times_A A'_2 \longrightarrow
B' =
(A'_1 \otimes_{A_1} B_1)_{h_1} \times_B
(A'_2 \otimes_{A_1} B_2)_{h_2}
$$
By
More on Algebra, Lemma \ref{more-algebra-lemma-module-over-fibre-product}
we have $A'_1 \otimes_{A'} B' = (A'_1 \otimes_{A_1} B_1)_{h_1}$
and $A'_2 \otimes_{A'} B' = (A'_2 \otimes_{A_2} B_2)_{h_2}$. In particular
the fibres of the morphism $Y' = \Spec(B') \to \Spec(A') = X'$ are
open subschemes of base changes of the fibres of the maps $Y_i \to X_i$,
hence finite and local complete intersections
(see Lemma \ref{lemma-syntomic-quasi-finite}).
By More on Algebra, Lemma
\ref{more-algebra-lemma-properties-algebras-over-fibre-product}
the ring map $A' \to B'$ is flat and of finite presentation.
Thus by Lemma \ref{lemma-syntomic-quasi-finite} part (3)
we conclude that $Y' \to X'$ is locally quasi-finite and syntomic.
Thus we obtain the commutative diagram
$$
\xymatrix{
& & Y/X \ar[ld] \ar@/_2pc/[lldd]_{b_1/a_1} \ar[rd] \ar@/^2pc/[rrdd]^{b_2/a_2} \\
& Y'_1/X'_1 \ar[rd] \ar[ld] & & Y'_2/X'_2 \ar[ld] \ar[rd] \\
Y_1/X_1 & & Y'/X' & & Y_2/X_2
}
$$
where $Y'_i/X'_i$ corresponds to $A'_i \to (A'_i \otimes_{A_i} B_i)_{h_i}$.

\medskip\noindent
In this paragraph, we fix the issue that $A'$ and $A'_i$ may not be Noetherian
by a limit argument; we suggest the reader skip this.
We can find a finitely generated $\mathbf{Z}$-subalgebra
$A'' \subset A'$ and a quasi-finite and syntomic ring map
$A'' \to B''$ such that $B' = A' \otimes_{A''} B''$.
This follows from Limits, Lemmas
\ref{limits-lemma-descend-finite-presentation},
\ref{limits-lemma-descend-syntomic}, and
\ref{limits-lemma-descend-quasi-finite}.
Then the ring maps $A'' \to A'_1 = A_1[x_\alpha]$ will factor through a finite
polynomial subalgebra $A_i[x_1, \ldots, x_n$ and the ring map
$A'' \to A'_2 = A_2[y_\beta]$ will factor through $A_2[y_1, \ldots, y_m]$.
After increasing the sets of variables, we may also assume that
$h_1 \in A_1[x_1, \ldots, x_n] \otimes_{A_1} B_1$ and
$h_2 \in A_2[y_1, \ldots, y_m] \otimes_{A_2} B_2$.
Increasing the sets of variables one more time, we can assume
that the maps $B'' \to (A'_i \otimes_{A_i} B_i)_{h_i}$
factor through $(A_1[x_1, \ldots, x_n] \otimes_{A_1} B_1)_{h_1}$
and $(A_2[y_1, \ldots, y_m] \otimes_{A_2} B_2)_{h_2}$.
Replacing $Y'$, $X'$, $Y'_1$, $X'_1$, $Y'_2$, $X'_2$ by the spectra of
$B''$, $A''$, $(A_1[x_1, \ldots, x_n] \otimes_{A_1} B_1)_{h_1}$,
$A_1[x_1, \ldots, x_n]$, $(A_2[y_1, \ldots, y_m] \otimes_{A_2} B_2)_{h_2}$,
$A_2[y_1, \ldots, y_m]$, we see that we get a diagram as above
where all schemes are Noetherian.

\medskip\noindent
Note that $Y'_i/X'_i$ is an object of $\mathcal{C}_{nice}$ as it
obtained by taking an open subscheme of the product of an affine space
with $Y_i/X_i$. In particular, the pullback of
$c_{Y_i/X_i}$ via $(b_i/a_i)$ is the same as the pullback of
$c_{Y'_i/X'_i}$ via $Y/X \to Y'_i/X'_i$.
Now we would be done if $Y'/X'$ is an object of $\mathcal{C}_{nice}$
(but this is likely not the case). Namely, then pulling back $c_{Y'/X'}$
around the two sides of the square, we would obtain the desired conclusion.
To get around the problem that $Y'/X'$ is not in $\mathcal{C}_{nice}$
we note the arguments above show that, after possibly shrinking all
of the schemes $X, Y, X'_1, Y'_1, X'_2, Y'_2, X', Y'$ we can find some
$n, d \geq 1$, and extend the diagram like so:
$$
\xymatrix{
& Y/X \ar[ld] \ar[rd] \\
Y'_1/X'_1 \ar[rd] & & Y'_2/X'_2 \ar[ld] \\
& Y'/X' \ar[d] \\
& Y_{n, d}/X_{n, d}
}
$$
and then we can use the already given argument by pulling
back from $c_{Y_{n, d}/X_{n, d}}$. This finishes the proof.
\end{proof}

\end{document}
